\documentclass[12pt,a4paper]{article}
\usepackage[margin=1in]{geometry}
\usepackage{times}
\usepackage[hidelinks]{hyperref}
\usepackage{booktabs}
\usepackage{array}
\usepackage{amsmath}
\usepackage{amssymb}
\usepackage{graphicx}
\usepackage{url}
\usepackage{microtype}
\usepackage{fancyhdr}
\usepackage{setspace}
\usepackage{enumitem}
\usepackage{multirow}
\usepackage{longtable}
\usepackage{parskip}
\usepackage[T1]{fontenc}
\usepackage{lmodern}

\pagestyle{fancy}
\fancyhf{}
\fancyhead[L]{\small Literature Review Report}
\fancyhead[R]{\small Big Data Analytics -- CS 5312}
\fancyfoot[C]{\thepage}
\renewcommand{\headrulewidth}{0.4pt}

\setlength{\parindent}{0pt}
\setlength{\parskip}{6pt}

\begin{document}

% ─── Title Block ──────────────────────────────────────────────────────────────
\begin{center}
{\LARGE \textbf{Literature Review: Auditing Geographic and Cultural Bias\\[0.4em]
in Clinical Large Language Models}}\\[1.8em]
{\large \textbf{Group Members:}}\\[0.6em]
Shawal Latif\\
Usmar Haider\\
Taimoor Karim\\
Abdul Moeed Irshad\\
Syed Muhammad Mujtaba\\[1.2em]
\textbf{Course:} Big Data Analytics -- CS 5312 \qquad
\textbf{Semester:} Spring 2026\\
\textbf{Instructor:} Dr.\ Imdadullah Khan\\
\textbf{Institution:} Lahore University of Management Sciences (LUMS)
\end{center}

\vspace{0.8em}\hrule\vspace{0.8em}

% ─── Abstract ─────────────────────────────────────────────────────────────────
\begin{abstract}
This literature review surveys fourteen papers at the intersection of clinical
artificial intelligence, algorithmic fairness, and large language model (LLM)
evaluation, collectively motivating a systematic audit of geographic and cultural
bias in clinical LLMs. The reviewed work is organised into four thematic
categories: (1)~racial and sociodemographic bias in clinical AI systems;
(2)~geographic and cultural bias in language models; (3)~clinical LLM evaluation
frameworks and benchmarks; and (4)~bias-detection methodologies based on
perturbation and counterfactual analysis. The literature shows that LLMs
consistently produce disparate clinical recommendations driven by non-clinical
identity cues---patient name, race, dialect, and geographic origin---even when
the underlying medical presentation is held constant. Crucially, however, every
reviewed clinical bias study operates within Global-North demographic categories;
the Global North/South axis has never been evaluated as a systematic perturbation
dimension. This review identifies that gap as the central motivation for our
project, which introduces the Geographic Disparity Index (GDI) as a novel
fairness metric and applies semantic-preserving geographic perturbation across
seven LLMs and four clinical datasets to quantify geographic bias in
AI-assisted clinical decision-making.
\end{abstract}

% ─── 1. Introduction ──────────────────────────────────────────────────────────
\section{Introduction}

Large Language Models (LLMs) are increasingly being deployed or considered for
deployment in clinical decision-support roles, particularly in regions facing
acute physician shortages~\cite{petterson2012}. In the Global South---encompassing
South Asia, Sub-Saharan Africa, Southeast Asia, the Middle East, and Latin
America---over six billion people increasingly interact with or are targeted by
AI-assisted clinical tools, presenting both enormous promise and substantial risk.
The promise lies in democratising access to medical guidance; the risk lies in
whether models trained predominantly on English-language Western biomedical text
encode systematic biases that disadvantage patients from underrepresented
geographic and cultural backgrounds.

This research area sits at the confluence of \emph{big data analytics},
\emph{natural language processing}, and \emph{health equity}. From a data
perspective the scale of the problem is immense: modern LLMs are trained on
corpora containing billions of tokens, and even small systematic asymmetries in
training data---for example, an over-representation of patients named ``David
from New York'' relative to ``Ali from Karachi''---can be amplified at inference
time across millions of clinical interactions globally. Detecting and quantifying
such biases requires large-scale experimental designs involving tens of thousands
of perturbation trials across multiple models, datasets, and demographic
dimensions---precisely the kind of task for which big-data pipelines are
indispensable.

The specific problem we address is: \emph{Do clinical LLMs recommend different
levels of care for patients with identical medical presentations when the
patient's geographic or cultural identity is varied through name and location
signals?} This is a health-equity question with direct societal consequences:
a model that systematically under-refers Global South patients for specialist
evaluation or diagnostic workup reproduces structural inequities at algorithmic
scale with an appearance of neutrality.

The objectives of this literature review are: (1)~to survey existing evidence
for demographic and geographic bias in clinical AI systems; (2)~to review
methodologies for bias detection in LLMs, particularly perturbation-based
approaches; (3)~to assess clinical LLM evaluation frameworks relevant to our
experimental design; and (4)~to identify the research gaps our project is
positioned to fill. The review is organised as follows: Section~2 provides
background and formal definitions; Section~3 surveys the literature in four
thematic categories; Section~4 provides comparative analysis; Section~5
identifies research gaps; Section~6 articulates the connection to our project;
and Section~7 concludes.

% ─── 2. Background ────────────────────────────────────────────────────────────
\section{Background and Preliminaries}

\subsection{Key Definitions}

\textbf{Large Language Model (LLM).}
A transformer-based neural network trained on large text corpora via
next-token prediction~\cite{thirunavukarasu2023}. In clinical contexts LLMs
are prompted with patient descriptions, symptoms, and queries to generate
free-text management recommendations.

\textbf{Algorithmic Bias.}
A systematic and unjustified deviation in model outputs for members of
different demographic groups~\cite{blodgett2020}. Bias may arise from training
data imbalance, proxy variables, or the implicit encoding of societal
disparities.

\textbf{Demographic Perturbation.}
An evaluation technique in which non-clinical identity signals (e.g., patient
name, city of origin) are varied while clinical content is held constant,
to test whether model outputs change~\cite{gourabathina2025,wan2023}.

\textbf{Semantic-Preserving Perturbation.}
A specific variant in which all medically relevant tokens---symptoms, history,
medications, examination findings---are held strictly constant; only the
identity-signalling tokens are modified.

\textbf{Global North / Global South.}
A socioeconomic and geopolitical distinction between historically industrialised
nations (US, UK, Europe, Australia) and developing nations (South Asia,
Sub-Saharan Africa, Southeast Asia, MENA, Latin America).

\textbf{Clinical Decision Support (CDS).}
AI- or algorithm-assisted guidance for medical practitioners or patients about
diagnosis, treatment, or care escalation.

\subsection{Foundational Techniques}

\textbf{Transformer Architecture and Attention.}
LLMs encode contextual information through multi-head self-attention. The
attention between position $i$ and position $j$ is:
\begin{equation}
    \operatorname{Attention}(Q,K,V)
    = \operatorname{softmax}\!\!\left(\frac{QK^{\!\top}}{\sqrt{d_k}}\right)V
\end{equation}
where $Q,K,V$ are query, key, and value matrices and $d_k$ is the key
dimension. Name tokens and geographic tokens may attract disproportionate
attention, causing identity signals to modulate clinical
reasoning~\cite{hofmann2024}.

\textbf{LLM-as-Annotator.}
Given the scale of output annotation required (${\sim}75{,}000$ model
responses), human annotation is infeasible. Following Zheng et
al.~\cite{zheng2023} we use a secondary LLM to extract structured binary
judgements from model outputs. For a model response $r$, the annotator
extracts:
\begin{align}
    y_{\text{MANAGE}}   &= \mathbf{1}[\text{response recommends home self-management}]\\
    y_{\text{VISIT}}    &= \mathbf{1}[\text{response recommends clinic or ED visit}]\\
    y_{\text{RESOURCE}} &= \mathbf{1}[\text{response recommends diagnostic resource}]
\end{align}

\subsection{Formal Problem Statement and Novel Metrics}

Let $\mathcal{C}$ be a set of clinical cases with ground-truth clinician
annotations. For each case $c\in\mathcal{C}$, let $x_c^{(0)}$ denote the
baseline (Global North) prompt and $x_c^{(k)}$ the $k$-th geographically
perturbed variant. Let $f_\theta$ denote an LLM with parameters $\theta$.

\textbf{Treatment Shift Rate (TSR):}
\begin{equation}
    \mathrm{TSR}
    = \frac{1}{|\mathcal{C}|}
      \sum_{c\in\mathcal{C}}
      \mathbf{1}\!\left[f_\theta(x_c^{(k)})\neq f_\theta(x_c^{(0)})\right]
\end{equation}

\textbf{Reduced Care Rate (RCR):}
\begin{equation}
    \mathrm{RCR}
    = \frac{\bigl|\{c : \text{shift is care-reducing}\}\bigr|}{|\mathcal{C}|}
\end{equation}

\textbf{Reduced Care Error Rate (RCER):}
\begin{equation}
    \mathrm{RCER}
    = \frac{\bigl|\{c : \text{shift is care-reducing}
                        \;\wedge\; \text{contradicts clinician label}\}\bigr|}
           {\bigl|\{c : \text{shift is care-reducing}\}\bigr|}
\end{equation}

\textbf{Geographic Disparity Index (GDI)} -- our novel composite metric:
\begin{equation}
    \mathrm{GDI}
    = \frac{1}{|\mathcal{R}|}
      \sum_{r\in\mathcal{R}} w_r\,
      \bigl(\mathrm{RCER}^{(\mathrm{South})}_r
            - \mathrm{RCER}^{(\mathrm{North})}_r\bigr)
\end{equation}
where $\mathcal{R}$ is the set of geographic regions, $w_r$ is a
population-proportional weight, and $\mathrm{RCER}^{(\cdot)}_r$ is the RCER
under Global South and Global North perturbations respectively.
$\mathrm{GDI}>0$ with $p<0.005$ (Bonferroni-corrected) indicates systematic
geographic bias.

% ─── 3. Literature Survey ─────────────────────────────────────────────────────
\section{Literature Survey}

\subsection{Category 1: Racial and Sociodemographic Bias in Clinical AI Systems}

Clinical AI systems have repeatedly been shown to encode and amplify
sociodemographic disparities. This section surveys four foundational papers
establishing the scope and severity of this problem.

\subsubsection*{Obermeyer et al.\ \cite{obermeyer2019}: Dissecting Racial Bias
in a Health Algorithm}

\textbf{Problem addressed.}
A widely deployed commercial health-risk algorithm was suspected of
systematically under-assigning risk scores to Black patients relative to equally
ill White patients, leading to fewer referrals for supplemental care programmes.

\textbf{Approach.}
The authors reverse-engineered the algorithm by analysing 43,539 primary-care
patients, discovering that it used \emph{healthcare costs} as a proxy for
\emph{health needs}---a proxy corrupted by structural racism, since Black patients
historically spend less on healthcare even when equally sick due to systemic
access barriers.

\textbf{Key contributions.}
The paper provides canonical evidence that algorithmic health proxies can
reproduce racial disparities without containing explicit racial features. It
establishes the auditing methodology of comparing algorithm outputs across
matched cohorts that differ only in race.

\textbf{Dataset.}
43,539 patients from a large academic medical centre; outcome labels from EHR
records.

\textbf{Results.}
At a given risk score, Black patients were found to be significantly sicker than
White patients. The algorithm effectively assigned lower risk to Black patients
despite higher disease burden, reducing Black enrolment in care management
programmes by 50\%.

\textbf{Limitations.}
Focuses on a single proprietary algorithm at a single institution. The bias
mechanism (cost as proxy) may not generalise to LLM-based systems that operate
on natural language.

\subsubsection*{Zack et al.\ \cite{zack2024}: GPT-4 Perpetuates Racial and
Gender Biases in Healthcare}

\textbf{Problem addressed.}
Whether GPT-4 reproduces documented racial and gender biases encoded in medical
education when performing clinical vignette tasks.

\textbf{Approach.}
The authors tested GPT-4 on four clinical vignettes known to contain erroneous
race- and gender-based clinical recommendations (e.g., race-adjusted kidney
function estimates), evaluating whether GPT-4 perpetuated these biases in
free-text responses.

\textbf{Key contributions.}
Direct evidence that GPT-4 reproduces race-based medical myths (e.g.,
spirometry race adjustment, differential creatinine correction for Black
patients) encoded in historical medical texts. Establishes LLMs as potential
vectors for disseminating debunked race-based pseudoscience.

\textbf{Dataset.}
4 clinical vignettes from established medical bias literature; 5 clinical tasks.

\textbf{Results.}
GPT-4 incorporated race as a clinical variable in 65\% of vignettes where it
should not, reproducing documented racial biases in pain management, kidney
function, and obstetric risk assessment.

\textbf{Limitations.}
Limited to 4 vignettes and a single model. Does not apply counterfactual
perturbation systematically. Exclusively U.S.-centric racial categories.

\subsubsection*{Omiye et al.\ \cite{omiye2023}: Large Language Models Propagate
Race-Based Medicine}

\textbf{Problem addressed.}
Whether four widely used LLMs (GPT-4, GPT-3.5, Bard, Llama-2) perpetuate
race-based medical practices that lack biological justification.

\textbf{Approach.}
The authors identified ten documented examples of debunked race-based medicine
from the medical literature and queried each model on these topics, then
evaluated response accuracy and racial bias through qualitative content analysis.

\textbf{Key contributions.}
First systematic multi-model comparison of race-based medical bias in LLMs.
Demonstrates that all tested models perpetuate at least some race-based medical
falsehoods and introduces a content-analysis framework for evaluating
race-based LLM outputs.

\textbf{Dataset.}
10 clinical topics $\times$ 4 models $\times$ 5 repetitions.

\textbf{Results.}
8 of 10 topics elicited at least one model perpetuating false race-based claims.
GPT-3.5 performed worst; Llama-2 and GPT-4 showed better but not perfect
performance.

\textbf{Limitations.}
Binary evaluation (propagates bias yes/no) does not measure degree or clinical
impact. Does not evaluate geographic or cultural origin as a perturbation axis.

\subsubsection*{Omar et al.\ \cite{omar2025}: Sociodemographic Biases in Medical
Decision-Making by LLMs}

\textbf{Problem addressed.}
Whether LLMs produce different clinical management recommendations for emergency
department (ED) patients based on sociodemographic characteristics when clinical
presentations are held constant.

\textbf{Approach.}
1,000 ED cases were presented to 9 LLMs (GPT-4, Claude, Llama, Gemini), each
in 32 sociodemographic variants (race, gender, insurance, housing, LGBTQIA+
status, BMI), yielding 1.7 million total outputs. Differences in diagnostic
workup, specialist referral, and treatment intensity were measured.

\textbf{Key contributions.}
Largest-scale study to date on sociodemographic bias in clinical LLM
decision-making. Finds that Black, unhoused, and LGBTQIA+ patients are directed
toward invasive interventions or urgent care with no clinical justification.

\textbf{Dataset.}
1,000 real ED cases; 32 demographic variants; 9 LLMs; 1.7 million outputs.

\textbf{Results.}
Significant ($p<0.001$) disparities across all tested sociodemographic
dimensions. GPT-4 showed the most consistent disparities; Llama-3 and
Gemini-Pro showed smaller but statistically significant effects.

\textbf{Limitations.}
Exclusively U.S.-based sociodemographic categories. Does not test geographic
origin (Global North vs.\ Global South), multilingual models, or
biomedical-specialised models. Perturbation limited to explicitly declared
demographic attributes rather than name-based inference.

% ──────────────────────────────────────────────────────────────────────────────
\subsection{Category 2: Geographic and Cultural Bias in Language Models}

While clinical AI literature has focused on within-U.S.\ racial disparities,
a parallel body of NLP work has established that language models encode
geographic and cultural biases intrinsically. This section reviews four papers
establishing that LLMs are systematically underrepresentative of the Global South.

\subsubsection*{Gourabathina et al.\ \cite{gourabathina2025}: The Medium is the
Message}

\textbf{Problem addressed.}
Whether non-clinical information in patient messages---patient name, writing
style, and perceived socioeconomic status---influences LLM-generated clinical
responses in ways that parallel documented human provider bias.

\textbf{Approach.}
The authors applied a semantic-preserving perturbation framework to a clinical
QA dataset, varying patient names (Black-associated vs.\ White-associated names
derived from U.S.\ Census data) and writing style (formal vs.\ informal) while
holding all clinical content constant. GPT-4 was the evaluated model; a second
GPT-4 instance served as the LLM-as-annotator.

\textbf{Key contributions.}
Establishes the direct methodological template for our project. Introduces the
three-question annotation schema (MANAGE / VISIT / RESOURCE) and the
semantic-preserving perturbation approach we adopt and extend to global
geographies.

\textbf{Dataset.}
OncQA (cancer patient-clinician messages) and AskDocs (Reddit medical queries);
$\sim$208 total cases with name perturbation.

\textbf{Results.}
Black-name patients received significantly more self-management recommendations
(Reduced Care Rate: $+12.3\%$ vs.\ White-name baseline) and fewer specialist
referrals. Informal writing style amplified the effect.

\textbf{Limitations.}
Restricted to U.S.-centric Black/White name pairs. Tests only GPT-4. Does not
extend to Global South geographies or multilingual models.

\subsubsection*{Hofmann et al.\ \cite{hofmann2024}: AI Generates Covertly Racist
Decisions Based on Dialect}

\textbf{Problem addressed.}
Whether LLMs make covertly racist decisions when input text is written in
African American English (AAE) versus Standard American English (SAE), even
when the model does not explicitly recognise the dialect.

\textbf{Approach.}
A dialect classifier was used to systematically vary input text between AAE and
SAE across six decision-making tasks, including criminal sentencing and
employment screening. Multiple LLMs were evaluated on the resulting disparate
output distributions.

\textbf{Key contributions.}
Demonstrates that linguistic form---not explicit race labels---can trigger
discriminatory model behaviour, bypassing explicit racial fairness guardrails.
Introduces the \emph{covert racism} framing: models that deny racial bias while
exhibiting it through linguistic proxies.

\textbf{Dataset.}
Dialect corpus with 200+ texts per condition; 6 decision tasks; multiple LLMs.

\textbf{Results.}
LLMs exhibited significantly more negative stereotyping for AAE texts.
Effects persisted when models were explicitly instructed to ignore race.

\textbf{Limitations.}
Focuses on U.S.\ dialect, not cross-national linguistic variation. Does not
examine clinical domains or geographic origin.

\subsubsection*{Manvi et al.\ \cite{manvi2024}: Large Language Models are
Geographically Biased}

\textbf{Problem addressed.}
Whether LLMs have systematic geographic biases in knowledge and prediction,
over-representing wealthy Western regions.

\textbf{Approach.}
The geographic knowledge of GPT-4, Gemini, and Llama-2 was probed using
geographic prediction tasks (e.g., estimating crop yields, poverty rates, and
satellite-imagery features from text descriptions). Prediction errors were
correlated with GDP, population density, and Global North/South classification.

\textbf{Key contributions.}
First systematic study of geographic bias in LLMs across multiple dimensions.
Shows that prediction errors correlate strongly with GDP ($r=-0.62$), meaning
the Global South is systematically underserved by LLM geographic knowledge.

\textbf{Dataset.}
Worldwide geographic prediction tasks across 170+ countries; GDP, HDI, and
population data from the World Bank and UN.

\textbf{Results.}
LLMs significantly over-perform on Global North regions and under-perform on
Global South regions. The effect is not correctable by simple recalibration.

\textbf{Limitations.}
Tests geographic prediction tasks (agriculture, economics) rather than clinical
reasoning. Does not directly evaluate whether geographic bias transfers to
healthcare decision-making.

\subsubsection*{Faisal and Anastasopoulos \cite{faisal2022}: Geographic and
Geopolitical Biases of Language Models}

\textbf{Problem addressed.}
Whether pretrained language models encode geographic and geopolitical biases
in their representations, beyond surface-level language differences.

\textbf{Approach.}
A Geographic Representation Probing Framework was developed using probing
classifiers to test whether LM representations encode geographic location,
geopolitical alignment, and wealth/development status, evaluated on mBERT,
XLM-R, and GPT-2 across 193 countries.

\textbf{Key contributions.}
Establishes that geographic and geopolitical biases are intrinsic to LM
representations and amplified at inference time. Models cluster countries along
geopolitical rather than geographic lines, reflecting training-data political
framing.

\textbf{Dataset.}
Wikipedia articles and news corpora for 193 countries; geopolitical alignment
labels from UN voting records.

\textbf{Results.}
Probing classifiers achieve 85\% accuracy at predicting geopolitical bloc from
LM representations. English-trained models show stronger Western bias.

\textbf{Limitations.}
Focuses on representational bias rather than downstream task performance. Does
not evaluate clinical or health-related tasks.

% ──────────────────────────────────────────────────────────────────────────────
\subsection{Category 3: Clinical LLM Evaluation Frameworks and Benchmarks}

Effective clinical LLM auditing requires robust evaluation frameworks. This
section surveys four papers that have developed methodologies for assessing LLM
performance and safety in clinical settings.

\subsubsection*{Thirunavukarasu et al.\ \cite{thirunavukarasu2023}: Large
Language Models in Medicine}

\textbf{Problem addressed.}
How do LLMs perform across the spectrum of clinical tasks, and what are the
key risks and evaluation requirements for equitable clinical deployment?

\textbf{Approach.}
A comprehensive narrative review of LLM performance across clinical tasks
(diagnosis, treatment recommendation, documentation, patient communication)
with meta-analysis of published benchmarks (MedQA, USMLE, ClinicalBench).

\textbf{Key contributions.}
Establishes the canonical taxonomy of clinical LLM capabilities and failure
modes. Identifies training-data bias as the primary risk to equitable clinical
deployment. Calls for standardised evaluation frameworks using diverse patient
populations.

\textbf{Dataset.}
Review of 110+ papers; benchmarks include USMLE-MedQA, MedMCQA, PubMedQA.

\textbf{Results.}
GPT-4 achieves passing scores on USMLE ($>70\%$) but performance degrades
significantly for non-English clinical scenarios and diseases prevalent in the
Global South (e.g., malaria, dengue, sickle-cell anaemia).

\textbf{Limitations.}
Does not systematically evaluate geographic or cultural bias. Reviewed
benchmarks are predominantly derived from U.S.\ and European clinical contexts.

\subsubsection*{Hager et al.\ \cite{hager2024}: Evaluation and Mitigation of LLM
Limitations in Clinical Decision-Making}

\textbf{Problem addressed.}
How to systematically identify LLM failure modes in clinical decision-making,
and whether mitigation strategies (chain-of-thought prompting, retrieval
augmentation) address identified limitations.

\textbf{Approach.}
A structured evaluation battery of 300 expert-curated clinical decision tasks
was applied to GPT-4, Claude-2, and Llama-2, evaluating both accuracy and
calibration. Mitigation strategies were then tested on identified failure modes.

\textbf{Key contributions.}
Introduces the first framework specifically designed to evaluate LLM
\emph{failure modes} in clinical settings, distinguishing hallucination,
knowledge gaps, and bias as distinct categories. Shows that chain-of-thought
prompting improves accuracy but does not reduce demographic bias.

\textbf{Dataset.}
300 expert-curated clinical decision tasks; validated by 12 clinician annotators.

\textbf{Results.}
All tested LLMs showed systematic overconfidence. Retrieval-augmented generation
reduced hallucination but had no significant effect on demographic bias.

\textbf{Limitations.}
Clinical tasks designed by U.S.\ and European clinicians. No systematic
geographic perturbation.

\subsubsection*{Pfohl et al.\ \cite{pfohl2024}: A Toolbox for Surfacing Health
Equity Harms in LLMs}

\textbf{Problem addressed.}
How to systematically identify and measure health equity harms produced by LLMs
in a way that is reproducible across research groups.

\textbf{Approach.}
The authors developed an evaluation toolbox comprising: (a)~457 clinical
scenarios with paired patient descriptions varying by race, gender, age, and
insurance; (b)~a structured annotation rubric; and (c)~reproducible disparity
metrics.

\textbf{Key contributions.}
First standardised toolbox for health-equity evaluation of clinical LLMs.
Demonstrates that GPT-4, Bard, and Llama-2-70B all exhibit measurable health
equity harms across multiple demographic dimensions.

\textbf{Dataset.}
457 clinical scenarios $\times$ 8 demographic variants; validated by 6 clinician
annotators.

\textbf{Results.}
All three tested LLMs exhibited significant demographic disparities in at least
3 of 6 clinical task categories. Insurance status and race produced the largest
disparities.

\textbf{Limitations.}
Geographic origin is not a tested dimension. All scenarios are U.S.-centric.

\subsubsection*{Poulain et al.\ \cite{poulain2024}: Bias Patterns in LLMs for
Clinical Decision Support}

\textbf{Problem addressed.}
Whether bias in clinical LLMs is model-specific or systematic across model
families, and which protected attributes drive the largest disparities.

\textbf{Approach.}
8 LLMs (GPT-4, Claude-2, Llama-2-70B, Falcon-40B, and others) were evaluated
across 3 clinical QA datasets using standardised patient vignettes. Demographic
attributes were varied via template substitution; output distributions were
compared using Kullback--Leibler divergence.

\textbf{Key contributions.}
Shows that clinical LLM bias is not model-specific---all tested models exhibit
significant disparities. Demonstrates that biomedical fine-tuned models show
\emph{larger} disparities than general-purpose models ($+23\%$ average), suggesting
fine-tuning on clinical corpora amplifies rather than reduces bias.

\textbf{Dataset.}
MedQA-USMLE, ClinicalNLP, MedPaLM evaluation sets; 3 demographic attributes
$\times$ 5 values.

\textbf{Results.}
Race and socioeconomic status produced the largest disparities
($\mathrm{KL}=0.31$ and $0.28$ respectively). Fine-tuned models showed $23\%$
larger average disparities than their base counterparts.

\textbf{Limitations.}
Does not test geographic origin or non-English names. All demographic signals
inserted via explicit attribute statements rather than implicit name-based
inference.

% ──────────────────────────────────────────────────────────────────────────────
\subsection{Category 4: Bias Detection Methodologies: Perturbation and
Counterfactual Analysis}

The methodological foundation of demographic bias auditing in NLP relies on
perturbation and counterfactual techniques. This section surveys three papers
that have shaped the methodological landscape most directly relevant to our work.

\subsubsection*{Blodgett et al.\ \cite{blodgett2020}: Language Technology is Power}

\textbf{Problem addressed.}
Why is NLP bias research methodologically inconsistent, and what normative
and empirical framework should guide it?

\textbf{Approach.}
A critical survey of 146 NLP papers on ``bias'', analysing their motivation,
operationalisation, and claims. The authors argue that most papers conflate
technical and normative questions without specifying whose interests are
protected or what harms are prevented.

\textbf{Key contributions.}
Seminal work establishing normative foundations for NLP bias research.
Introduces the distinction between \emph{representational} bias (stereotyping
in model representations) and \emph{allocational} bias (differential
performance on downstream tasks). Argues for grounding bias research in
documented, specific harms to named communities.

\textbf{Dataset.}
146 NLP papers reviewed; no new experimental data.

\textbf{Results.}
Most surveyed papers do not specify which communities are harmed or how. Few
ground their bias definitions in documented real-world harm.

\textbf{Limitations.}
Methodological critique rather than experimental paper. Does not propose
specific new evaluation datasets or metrics.

\subsubsection*{Wan et al.\ \cite{wan2023}: Gender Biases in LLM-Generated
Reference Letters}

\textbf{Problem addressed.}
Whether LLMs encode gender bias when generating professional reference letters,
and whether name-based perturbation reliably detects such bias.

\textbf{Approach.}
Female-stereotyped names (Kelly, Jennifer) were replaced with male-stereotyped
names (Joseph, David) in prompts requesting LLM-generated reference letters.
Generated text was analysed for agentic vs.\ communal language using LIWC and
word-embedding analyses across five LLMs.

\textbf{Key contributions.}
Validates name-based perturbation as a reliable probe for demographic bias in
LLM outputs. Shows that LLMs systematically describe female-named subjects with
warmth/likeability language and male-named subjects with leadership/competence
language even when the underlying request is identical. Establishes the
methodological precedent we extend from gender to geographic identity.

\textbf{Dataset.}
100 prompts $\times$ 2 name conditions $\times$ 5 LLMs (ChatGPT, GPT-4, Alpaca,
Dolly, Vicuna).

\textbf{Results.}
Female-named subjects received 68\% more communal language and 41\% less
agentic language, consistently across all tested LLMs.

\textbf{Limitations.}
Binary gender proxy; does not address non-binary or intersectional identity.
Limited to U.S.-centric names. Does not extend to geographic or cultural
identity.

\subsubsection*{Parrish et al.\ \cite{parrish2022}: BBQ: A Bias Benchmark for
Question Answering}

\textbf{Problem addressed.}
How to systematically measure whether QA models rely on social stereotypes when
answering questions about people from different demographic groups.

\textbf{Approach.}
A 58,492-example benchmark of question--context--answer triplets was constructed
in which the correct answer requires \emph{not} relying on demographic
stereotypes. Each question has an ambiguous version (only stereotypes provide a
cue) and a disambiguated version (context provides an unambiguous answer).

\textbf{Key contributions.}
BBQ is the most widely used benchmark for measuring stereotype reliance in LLMs.
Covers 9 social categories including race/ethnicity, nationality, religion,
gender, and age. Provides a replicable methodology for constructing
counterfactual evaluation sets directly applicable to clinical scenario design.

\textbf{Dataset.}
58,492 examples; 9 social categories; 3 question formats.

\textbf{Results.}
All tested models showed above-baseline reliance on stereotypes in ambiguous
contexts. Race/ethnicity and nationality categories produced the largest biases.

\textbf{Limitations.}
Focused on QA classification rather than generative outputs. Nationality
categories do not specifically distinguish Global North/South. Does not test
LLMs in clinical contexts.

% ─── 4. Comparative Analysis ──────────────────────────────────────────────────
\section{Comparative Analysis}

\subsection{Summary Table}

\begin{table}[h!]
\centering
\caption{Comparison of Reviewed Approaches}
\label{tab:comparison}
\small
\setlength{\tabcolsep}{4pt}
\begin{tabular}{p{3.0cm}p{3.0cm}p{2.8cm}p{2.8cm}p{3.0cm}}
\toprule
\textbf{Paper} & \textbf{Technique} & \textbf{Dataset(s)} &
\textbf{Key Finding} & \textbf{Limitations} \\
\midrule
Obermeyer et al.~\cite{obermeyer2019} &
Retrospective cohort audit &
43,539 EHR patients &
50\% less Black enrolment in care mgmt &
Single algorithm; single institution \\
\addlinespace
Zack et al.~\cite{zack2024} &
Clinical vignette testing &
4 clinical vignettes &
GPT-4 reproduces race-based medical myths &
4 cases; single model \\
\addlinespace
Omiye et al.~\cite{omiye2023} &
Multi-model topic audit &
10 topics $\times$ 4 LLMs &
All models perpetuate racial medical claims &
Binary evaluation only \\
\addlinespace
Omar et al.~\cite{omar2025} &
Large-scale demographic perturbation &
1,000 ED cases $\times$ 32 variants &
Disparities for Black/unhoused/LGBTQIA+ patients &
U.S.-centric; no geographic axis \\
\addlinespace
Gourabathina et al.~\cite{gourabathina2025} &
Semantic-preserving name perturbation &
OncQA, AskDocs &
Black-name patients get more self-management &
U.S.\ names only; single model \\
\addlinespace
Hofmann et al.~\cite{hofmann2024} &
Dialect substitution &
AAE/SAE corpus &
Covert racism via dialect proxy &
U.S.\ dialect only \\
\addlinespace
Manvi et al.~\cite{manvi2024} &
Geographic knowledge probing &
World Bank/UN, 170+ countries &
GDP-correlated geographic bias in LLMs &
Not clinical domain \\
\addlinespace
Faisal \& Anastasopoulos~\cite{faisal2022} &
Representation probing &
Wikipedia, 193 countries &
Geopolitical bias in LM representations &
Not downstream tasks \\
\addlinespace
Thirunavukarasu et al.~\cite{thirunavukarasu2023} &
Systematic literature review &
110+ papers; USMLE-MedQA &
LLMs degrade on non-Western inputs &
No geographic bias audit \\
\addlinespace
Hager et al.~\cite{hager2024} &
Structured failure evaluation &
300 expert clinical tasks &
Chain-of-thought improves accuracy, not bias &
U.S./EU-centric tasks \\
\addlinespace
Pfohl et al.~\cite{pfohl2024} &
Equity toolbox (457 scenarios) &
457 scenarios $\times$ 8 variants &
Insurance/race drive largest disparities &
No geographic axis \\
\addlinespace
Poulain et al.~\cite{poulain2024} &
KL-divergence perturbation &
MedQA, ClinicalNLP &
Fine-tuning amplifies bias by 23\% &
No geographic signals \\
\addlinespace
Wan et al.~\cite{wan2023} &
Name-based perturbation &
100 prompts $\times$ 5 LLMs &
Female names receive communal language &
Binary gender; U.S.-centric \\
\addlinespace
Parrish et al.~\cite{parrish2022} &
Counterfactual QA benchmark &
58,492 BBQ examples &
Models rely on stereotypes under ambiguity &
Classification only; no clinical \\
\bottomrule
\end{tabular}
\end{table}

\subsection{Critical Analysis}

\textbf{Common strengths.}
The surveyed literature demonstrates convergent validity: across diverse
methodologies---retrospective audits, template perturbation, representation
probing, and large-scale vignette studies---a consistent finding emerges that
AI systems produce systematically disparate outputs for different demographic
groups. The most rigorous studies~\cite{omar2025,pfohl2024} employ large-scale
designs with multiple demographic dimensions and multi-model comparisons,
yielding high statistical power. Methodological maturity is increasing: the
field is moving from single-model case studies~\cite{zack2024,omiye2023} toward
standardised, reproducible frameworks~\cite{pfohl2024,parrish2022}.

\textbf{Common limitations.}
Three limitations recur across virtually all reviewed papers:
\begin{enumerate}[leftmargin=2em]
    \item \textbf{U.S.-centric scope.} All 14 papers evaluate U.S.-based
    demographic categories exclusively or primarily. The Global North/South
    axis remains entirely unexplored in the clinical AI literature.
    \item \textbf{Limited model diversity.} Most studies test 1--4 models,
    typically GPT-4 plus one or two open-source alternatives. No study
    simultaneously evaluates biomedical-specialised models alongside frontier
    proprietary and multilingual models.
    \item \textbf{Explicit vs.\ implicit demographic signals.} Papers using
    explicit statements (``Patient is Black'') may underestimate real-world
    bias, since clinical LLMs in practice infer demographics from contextual
    cues (names, city names, writing style). Only Gourabathina et
    al.~\cite{gourabathina2025} and Wan et al.~\cite{wan2023} use implicit
    name-based signals, and neither extends to geographic identity.
\end{enumerate}

\textbf{Methodological trends.}
The field is evolving from small-scale case studies toward large-scale
systematic evaluations~\cite{omar2025}. LLM-as-annotator pipelines are
increasingly adopted for scalable evaluation~\cite{zheng2023}. The introduction
of composite fairness metrics (RCER, KL disparity) represents a maturation from
binary (biased/unbiased) to quantitative bias measurement---a trajectory our GDI
metric continues.

\textbf{Scalability.}
Studies with explicit demographic injection~\cite{pfohl2024,poulain2024} scale
readily to large-scale experiments but may underestimate bias that operates
through implicit inference. Studies with name-based perturbation~\cite{wan2023,
gourabathina2025} are more ecologically valid but require larger sample sizes to
achieve statistical power, since the signal is noisier. Our project resolves
this by combining both perturbation types (name-only, location-only, combined)
across 1,541 base cases.

% ─── 5. Research Gaps ─────────────────────────────────────────────────────────
\section{Research Gaps and Opportunities}

The surveyed literature reveals five critical gaps that our project is
positioned to address.

\textbf{Gap 1: Absence of the Global South as an evaluation axis.}
Every reviewed clinical AI bias study uses demographic categories internal to
the Global North (Black/White in the U.S., AAE dialect, insurance status). The
6+ billion people in the Global South who are increasingly using or being
considered for AI-assisted clinical tools are invisible in the current evaluation
literature. Our project fills this gap by introducing geographic origin as a
systematic perturbation dimension with culturally validated name banks across
five global regions.

\textbf{Gap 2: Single-model evaluation.}
With limited exceptions~\cite{omar2025,omiye2023}, reviewed papers test a
single model or narrow pairs. No study simultaneously evaluates proprietary
frontier models (GPT, Claude, Gemini), multilingual models (Qwen, Kimi), and
biomedical-specialised models (BioMedGemma) under identical conditions. This
prevents understanding of whether geographic bias is an architectural or
training-data phenomenon.

\textbf{Gap 3: Implicit name-based geographic inference.}
Existing studies either use explicit demographic statements or U.S.-centric name
pairs. No study tests whether LLMs implicitly infer geographic origin from
culturally coded names (Ali, Fatima, Kwame, Nguyen) and whether this inference
triggers disparate clinical recommendations.

\textbf{Gap 4: Absence of a geographic fairness metric.}
No standardised metric exists for measuring geographic bias severity in a way
comparable across models and datasets. Existing metrics (disparity ratio, KL
divergence) were designed for within-Western demographic comparisons. Our
Geographic Disparity Index (GDI) is the first scalar metric designed
specifically for cross-regional fairness comparison in clinical AI.

\textbf{Gap 5: Intersectional geographic--gender analysis.}
While intersectionality has been theorised in AI fairness~\cite{blodgett2020},
no clinical AI paper evaluates the interaction between geographic origin and
gender (e.g., whether Global South women face compounded disparities). Our
project includes this as a secondary analysis dimension.

% ─── 6. Relevance to Project ──────────────────────────────────────────────────
\section{Relevance to Your Project}

\textbf{Adopted techniques.}
\begin{itemize}[leftmargin=2em]
    \item \emph{Semantic-preserving perturbation}: adopted directly from
    Gourabathina et al.~\cite{gourabathina2025} and Wan et al.~\cite{wan2023},
    extended to geographic identity (name + city pairs across 5 global regions
    and $\sim$200 culturally validated names).
    \item \emph{LLM-as-annotator pipeline}: adopted from Zheng et
    al.~\cite{zheng2023}, using LLaMA-3-8B to extract binary MANAGE/VISIT/RESOURCE
    labels from $\sim$75{,}000 model responses at scale.
    \item \emph{Three-question annotation schema}: adopted directly from
    Gourabathina et al.~\cite{gourabathina2025}.
    \item \emph{Statistical framework}: ANOVA with Bonferroni correction and
    paired $t$-tests with FDR correction, following Pfohl et
    al.~\cite{pfohl2024} and Poulain et al.~\cite{poulain2024}.
\end{itemize}

\textbf{Baselines.}
Our primary baseline is the within-model, within-dataset comparison between
Global North (David from New York) and Global South (Ali from Karachi)
perturbation conditions. Secondary baselines draw on Gourabathina et
al.~\cite{gourabathina2025} (Black/White name disparity in U.S.\ context) and
Omar et al.~\cite{omar2025} (sociodemographic disparity in ED settings).

\textbf{Relevant datasets from the literature.}
\begin{itemize}[leftmargin=2em]
    \item \emph{OncQA}~\cite{chen2023}: used in Gourabathina et
    al.~\cite{gourabathina2025}; 61 filtered cancer patient-clinician message
    pairs in our study.
    \item \emph{AskDocs (r/AskaDocs)}: used in Gourabathina et
    al.~\cite{gourabathina2025}; 147 informal medical queries.
    \item \emph{USMLE-MedQA}~\cite{jin2020}: reviewed in Thirunavukarasu et
    al.~\cite{thirunavukarasu2023} and Hager et al.~\cite{hager2024};
    1{,}333 filtered vignettes across 12 specialties.
\end{itemize}

\textbf{Evaluation metrics.}
Our metrics extend those in the reviewed literature: TSR (analogous to
disparity ratio in Pfohl et al.~\cite{pfohl2024}), RCR and RCER (extending
Gourabathina et al.'s annotation schema), and the novel GDI (addressing Gap 4
above). Significance thresholds follow the Bonferroni-corrected $p<0.005$
standard used in Omar et al.~\cite{omar2025}.

\textbf{Extension beyond existing work.}
Our project extends the literature in three key dimensions: (1)~geographic scope
(Global North/South vs.\ within-U.S.\ racial categories); (2)~model diversity
(7 models spanning proprietary, open-source, multilingual, and
biomedical-specialised families); and (3)~a novel composite fairness metric
(GDI) enabling cross-model comparison not possible with existing tools.

% ─── 7. Conclusion ────────────────────────────────────────────────────────────
\section{Conclusion}

This literature review has surveyed 14 papers at the intersection of clinical AI
bias, LLM evaluation, and geographic/cultural fairness in NLP. The central
finding is that AI systems---from health-scoring algorithms to frontier LLMs---
consistently produce disparate clinical outcomes for different demographic groups
even when medical presentations are held constant. Methods have matured from
case studies to large-scale systematic evaluations, and LLM-as-annotator
pipelines have made annotation feasible at scale.

However, a critical gap persists: no existing study examines the Global
North/South axis as a perturbation dimension in clinical AI. All reviewed
clinical bias work operates within U.S.-centric racial categories, leaving 6+
billion Global South patients without evidence-based assessment of the AI systems
increasingly targeting them.

Our project addresses this gap through: (1)~semantic-preserving geographic
perturbation across five global regions using culturally validated name banks;
(2)~multi-model evaluation spanning proprietary, open-source, multilingual, and
biomedical-specialised LLMs; (3)~four validated clinical datasets covering
oncology, general medicine, and structured vignettes; and (4)~the Geographic
Disparity Index (GDI), a novel fairness metric enabling cross-model comparison.

Anticipated challenges, informed by the literature, include: name-to-region
inference ambiguity (partially addressed by combined name + city perturbation);
LLM-as-annotator reliability in clinical contexts (mitigated by human
validation on a random subsample following Zheng et al.~\cite{zheng2023}); and
the risk that name-based inferences conflate geographic origin with other
correlated variables (addressed through intersectional analysis of the geographic
$\times$ gender interaction). This review provides the empirical and
methodological foundation for a rigorous, scalable audit of geographic bias in
clinical LLMs---a study positioned to have direct implications for health equity
in the Global South.

% ─── References ───────────────────────────────────────────────────────────────
\newpage
\begin{thebibliography}{99}

\bibitem{gourabathina2025}
S.~Gourabathina, T.~Naous, and W.~Xu,
``The medium is the message: How non-clinical information shapes clinical
decisions in {LLMs},''
in \emph{Proc.\ ACM FAccT}, 2025.

\bibitem{zack2024}
T.~Zack, E.~Steinberg, R.~Nori, J.~Listgarten, J.~Naik,
C.~Valliani, V.~Sherr, T.~Pierson, B.~Mach, and N.~Shah,
``Assessing the potential of {GPT-4} to perpetuate racial and gender biases
in health care: A model evaluation study,''
\emph{The Lancet Digital Health}, vol.~6, no.~1, pp.\ e12--e22, 2024.

\bibitem{hofmann2024}
V.~Hofmann, J.~Kalluri, D.~Jurafsky, and S.~King,
``{AI} generates covertly racist decisions about people based on their
dialect,''
\emph{Nature}, vol.~633, pp.\ 147--154, 2024.

\bibitem{omiye2023}
J.~A.~Omiye, J.~C.~Lester, S.~Spichak, V.~Rotemberg, and R.~Daneshjou,
``Large language models propagate race-based medicine,''
\emph{npj Digital Medicine}, vol.~6, art.\ 195, 2023.

\bibitem{pfohl2024}
S.~Pfohl \textit{et al.},
``A toolbox for surfacing health equity harms and biases in large language
models,''
\emph{Nature Medicine}, vol.~30, pp.\ 3333--3348, 2024.

\bibitem{hager2024}
P.~Hager, F.~Jungmann, R.~Holland, K.~Bhatt, E.~Hubber, B.~Knauer,
F.~Jakab, N.~Bhatt, A.~Lottermoser, and D.~Rueckert,
``Evaluation and mitigation of the limitations of large language models in
clinical decision-making,''
\emph{Nature Medicine}, vol.~30, pp.\ 2613--2622, 2024.

\bibitem{zheng2023}
L.~Zheng, W.-L.~Chiang, Y.~Sheng, S.~Zhuang, Z.~Wu, Y.~Zhuang, Z.~Li,
Z.~Li, D.~Li, E.~P.~Xing, H.~Zhang, J.~E.~Gonzalez, and I.~Stoica,
``Judging {LLM}-as-a-judge with {MT-Bench} and chatbot arena,''
in \emph{Proc.\ NeurIPS}, vol.~36, 2023.

\bibitem{obermeyer2019}
Z.~Obermeyer, B.~Powers, C.~Vogeli, and S.~Mullainathan,
``Dissecting racial bias in an algorithm used to manage the health of
populations,''
\emph{Science}, vol.~366, no.~6464, pp.\ 447--453, 2019.

\bibitem{blodgett2020}
S.~L.~Blodgett, S.~Barocas, H.~Daumé~III, and H.~Wallach,
``Language (technology) is power: A critical survey of `bias' in {NLP},''
in \emph{Proc.\ ACL}, pp.\ 5454--5476, 2020.

\bibitem{wan2023}
Y.~Wan, G.~Pu, J.~Sun, A.~Garimella, K.-W.~Chang, and N.~Peng,
```Kelly is a warm person, Joseph is a role model': Gender biases in
{LLM}-generated reference letters,''
in \emph{Findings of EMNLP}, pp.\ 3730--3748, 2023.

\bibitem{manvi2024}
R.~Manvi, S.~Khanna, M.~Burke, D.~Lobell, and S.~Ermon,
``Large language models are geographically biased,''
\emph{arXiv preprint arXiv:2402.02680}, 2024.

\bibitem{faisal2022}
F.~Faisal and A.~Anastasopoulos,
``Geographic and geopolitical biases of language models,''
in \emph{Proc.\ Workshop on Multilingual Representation Learning (MRL)
at EMNLP}, pp.\ 115--129, 2023.

\bibitem{thirunavukarasu2023}
A.~J.~Thirunavukarasu, D.~S.~J.~Ting, K.~Elangovan, L.~Gutierrez,
T.~F.~Tan, and D.~S.~W.~Ting,
``Large language models in medicine,''
\emph{Nature Medicine}, vol.~29, pp.\ 1930--1940, 2023.

\bibitem{poulain2024}
R.~Poulain, H.~Fayyaz, and R.~Beheshti,
``Bias patterns in the application of {LLMs} for clinical decision support:
A comprehensive study,''
\emph{arXiv preprint arXiv:2404.15149}, 2024.

\bibitem{omar2025}
M.~Omar, S.~Soffer, R.~Agbareia, N.~L.~Bragazzi, D.~U.~Apakama,
C.~R.~Horowitz, A.~W.~Charney, R.~Freeman, B.~Kummer,
B.~S.~Glicksberg, G.~N.~Nadkarni, and E.~Klang,
``Sociodemographic biases in medical decision making by large language
models,''
\emph{Nature Medicine}, vol.~31, pp.\ 1873--1881, 2025.

\bibitem{parrish2022}
A.~Parrish, A.~Chen, N.~Nangia, V.~Padmakumar, J.~Phang, J.~Thompson,
P.~M.~Htut, and S.~Bowman,
``{BBQ}: A hand-built bias benchmark for question answering,''
in \emph{Findings of ACL}, pp.\ 2086--2105, 2022.

\bibitem{chen2023}
E.~Chen, D.~Aggarwal, J.~Lee, K.~Ash, B.~Chen, and N.~Shah,
``The impact of using an {AI} chatbot to respond to patient messages,''
\emph{arXiv preprint arXiv:2310.17703}, 2023.

\bibitem{jin2020}
D.~Jin, E.~Pan, N.~Oufattole, W.-H.~Weng, H.~Fang, and P.~Szolovits,
``What disease does this patient have? A large-scale open domain question
answering dataset from medical exams,''
\emph{arXiv preprint arXiv:2009.13081}, 2020.

\bibitem{petterson2012}
S.~M.~Petterson, W.~R.~Liaw, R.~L.~Phillips, D.~L.~Rabin, D.~S.~Meyers,
and A.~W.~Bazemore,
``Projecting {US} primary care physician workforce needs: 2010--2025,''
\emph{Annals of Family Medicine}, vol.~10, no.~6, pp.\ 503--509, 2012.

\end{thebibliography}

\end{document}
